% ---------------------------------------------------------------------
%  CHƯƠNG 3 — PHÂN TÍCH VÀ THIẾT KẾ HỆ THỐNG
% ---------------------------------------------------------------------
\chapter{Phân tích và thiết kế hệ thống}
\label{ch:thietke}

\begin{muctieu}
Phân tích chi tiết 21 chức năng bài toán, thiết kế kiến trúc hệ thống 3 tầng,
mô hình hóa dữ liệu bất biến và các sơ đồ lớp UML, sơ đồ tuần tự chuẩn mực.
\end{muctieu}

\section{Phân tích bài toán và phân rã 21 chức năng}

Hệ thống trò chơi đoán số nhiều người được phân rã thành 21 chức năng chuẩn
mực từ F01 đến F21, phân định rõ khối lượng công việc giữa hai thành viên
trong nhóm:
\begin{itemize}[leftmargin=1.6em, itemsep=2pt]
  \item \textbf{Huỳnh Tiến Đạt (Logic nghiệp vụ \& Dữ liệu Streams):} Đảm nhiệm
        11 mục cốt lõi gồm mô hình \code{Game}, \code{LuotDoan},
        \code{TrangThaiVan}, \code{GoiY}, bảng điểm \code{KhoBangDiem}, lịch
        sử \code{KhoLichSuDoan}, thống kê \code{ThongKeService}, chế độ đoán
        lần lượt \code{ArrayDeque} và hẹn giờ ván với \code{Future} (F02, F03,
        F04, F05, F12, F13, F14, F15, F16, F19, F20).
  \item \textbf{Trần Thanh Bình (Hạ tầng mạng TCP, Giao diện \& Đa phòng):}
        Đảm nhiệm 10 mục hạ tầng gồm máy chủ \code{GameServer}, giao diện
        người dùng \code{ClientUI}, bể luồng \code{ClientHandler}, sổ người
        chơi online \code{SoNguoiChoiOnline}, broadcast tin nhắn, mở ván chơi,
        hiển thị dữ liệu bảng điểm/lịch sử/thống kê, xử lý ngắt kết nối dọn
        dẹp tài nguyên và điều phối đa phòng chơi song song chống Deadlock
        (F01, F06, F07, F08, F09, F10, F11, F17, F18, F21).
\end{itemize}

\section{Thiết kế kiến trúc hệ thống 3 tầng}

Kiến trúc hệ thống tuân thủ nguyên lý tách biệt mối quan tâm (Separation of
Concerns):
\begin{enumerate}[leftmargin=1.6em, itemsep=2pt]
  \item \textbf{Tầng Giao diện \& Mạng (\code{ui}, \code{net}):} Tiếp nhận yêu
        cầu người dùng, xử lý các kết nối Socket TCP, chuyển đổi chuỗi văn bản
        thành các lệnh gọi nghiệp vụ.
  \item \textbf{Tầng Nghiệp vụ (\code{domain}, \code{service}):} Xử lý toàn bộ
        logic tính toán, so sánh số bí mật, chuyển đổi trạng thái ván, xếp
        hàng lần lượt và điều phối phòng chơi. Tầng này độc lập hoàn toàn với
        Socket giúp việc viết unit test cực kỳ đơn giản và mạnh mẽ.
  \item \textbf{Tầng Lưu trữ (\code{repo}):} Đóng gói các kho lưu trữ dữ liệu
        an toàn luồng trong bộ nhớ (in-memory repositories) với cơ chế bảo vệ
        sao chép phòng thủ (defensive copying).
\end{enumerate}

\begin{figure}[H]
\centering
\includegraphics[width=0.95\textwidth]{02-kien-truc-cac-tang.drawio.png}
\caption{Sơ đồ kiến trúc tổng thể 3 tầng của hệ thống.}
\label{fig:kientruc}
\end{figure}

\section{Thiết kế dữ liệu và các thực thể Domain}

\begin{table}[H]
\centering
\caption{Cấu trúc các thực thể dữ liệu nghiệp vụ chính của hệ thống.}
\label{tab:dulieu}
\begin{tabular*}{\textwidth}{@{\extracolsep{\fill}}>{\bfseries\color{daunavy}}l >{\color{daumaroon}}l p{8.4cm}}
\toprule
\rowcolor{gray100}\bfseries\color{daunavy}Thực thể & \bfseries\color{daumaroon}Kiểu dữ liệu & \bfseries\color{daunavy}Mô tả và ý nghĩa thiết kế \\
\midrule
\code{LuotDoan}     & \code{record}           & Bản ghi bất biến: tên người chơi, số đoán, thời điểm (LocalDateTime), kết quả gợi ý \\
\code{TrangThaiVan} & \code{enum}             & Trạng thái vòng đời ván chơi: WAITING, PLAYING, FINISHED \\
\code{GoiY}         & \code{enum}             & Kết quả so sánh số đoán: HIGHER, LOWER, CORRECT \\
\code{CheDoDoan}    & \code{enum}             & Chế độ đoán của ván chơi: DONG\_THOI, LAN\_LUOT \\
\code{KetQua}       & \code{sealed interface} & Mô hình kết quả thống nhất: \code{ThanhCong} hoặc \code{ThatBai} \\
\code{PhongChoi}    & \code{class}            & Thực thể phòng chơi độc lập: có ván, bảng điểm, lịch sử và khóa đồng bộ riêng \\
\bottomrule
\end{tabular*}
\end{table}

\section{Thiết kế sơ đồ lớp UML}

\begin{figure}[H]
\centering
\includegraphics[width=0.98\textwidth]{03-so-do-lop-uml.drawio.png}
\caption{Sơ đồ lớp UML chi tiết của hệ thống trò chơi đoán số.}
\label{fig:uml}
\end{figure}

\section{Thiết kế luồng xử lý và tương tác tuần tự}

\begin{figure}[H]
\centering
\includegraphics[width=0.95\textwidth]{04-so-do-tuan-tu.drawio.png}
\caption{Sơ đồ tuần tự tương tác mạng giữa Máy khách và Máy chủ khi đoán số.}
\label{fig:tuantu}
\end{figure}
